\subsection*{CU-04: Verificar el correo y activar la cuenta}
\addcontentsline{toc}{subsection}{CU-04: Verificar el correo y activar la cuenta}
\label{cu-04}

\begin{fichacu}{CU-04}{Verificar el correo y activar la cuenta}
  \crudFila{Responsable}{José Eduardo Prior Hernández}
  \crudFila{Actor principal}{Jugador}
  \crudFila{Actores de sistema}{Servidor de partidas, Base de datos}
  \crudFila{Prototipos}{— (ocurre fuera de la aplicación, en el cliente de correo)}
  \crudFila{Objetivo}{Llevar una cuenta recién creada de \texttt{PENDIENTE} a \texttt{ACTIVA}, comprobando que quien se dio de alta controla la dirección de correo que declaró. Es el único flujo que puede activar una cuenta (\mbox{RN-04} del \mbox{CU-02}).}
  \crudFila{Disparador}{El Jugador abre el enlace de verificación que recibió por correo tras el \mbox{CU-02}, o tras el reenvío del \mbox{FA-09} del \mbox{CU-01}.}
  \textbf{Precondiciones} & \cuCelda{\begin{cureglas}
      \item \textbf{\mbox{PRE-01}} Existe un \textit{Usuario} con \texttt{estado\_\allowbreak{}cuenta} en \texttt{PENDIENTE}.
      \item \textbf{\mbox{PRE-02}} Existe un \textit{TokenVerificacionCorreo} asociado a ese \textit{Usuario}, sin usar.
      \item \textbf{\mbox{PRE-03}} El Servidor de partidas está en ejecución y tiene conexión disponible con la Base de datos.
    \end{cureglas}} \\ \hline
  \textbf{Flujo normal} & \cuCelda{\begin{cuenum}[start=1]
      \item El Jugador abre el enlace de verificación desde su cliente de correo.
      \item El Servidor de partidas recibe la solicitud con el token incluido en el enlace.
      \item El Servidor de partidas deriva el resumen criptográfico del token recibido y recupera de la Base de datos el \textit{TokenVerificacionCorreo} cuyo \texttt{token\_\allowbreak{}hash} coincida. \textit{(Ver \mbox{EX-01})}
      \item El Servidor de partidas verifica que el token exista. \textit{(Ver \mbox{FA-01})}
      \item El Servidor de partidas verifica que el token no esté marcado como usado. \textit{(Ver \mbox{FA-02})}
      \item El Servidor de partidas verifica que no haya expirado, comparando \texttt{fecha\_\allowbreak{}expiracion} con la fecha y hora actuales (\mbox{RN-02}). \textit{(Ver \mbox{FA-03})}
      \item El Servidor de partidas recupera el \textit{Usuario} asociado y verifica que su \texttt{estado\_\allowbreak{}cuenta} siga siendo \texttt{PENDIENTE}. \textit{(Ver \mbox{FA-04}, \mbox{FA-05}, \mbox{FA-07})}
      \item El Servidor de partidas inicia una transacción sobre la Base de datos.
    \end{cuenum}} \\
   & \cuCelda{\begin{cuenum}[start=9]
      \item El Servidor de partidas cambia el \texttt{estado\_\allowbreak{}cuenta} del \textit{Usuario} a \texttt{ACTIVA} y registra \texttt{fecha\_\allowbreak{}verificacion}. \textit{(Ver \mbox{EX-01})}
      \item El Servidor de partidas marca el \textit{TokenVerificacionCorreo} como usado (\mbox{RN-01}).
      \item El Servidor de partidas invalida los demás \textit{TokenVerificacionCorreo} pendientes del \textit{Usuario}, si los hubiera (\mbox{RN-03}).
      \item El Servidor de partidas confirma la transacción. \textit{(Ver \mbox{EX-02})}
      \item El Servidor de partidas registra la activación en la \textit{BitacoraAcceso} con el resultado \texttt{CUENTA\_\allowbreak{}VERIFICADA}.
      \item El Servidor de partidas muestra la página de confirmación en el \texttt{idioma\_\allowbreak{}preferido} del \textit{Usuario}, indicando que ya puede iniciar sesión.
      \item Fin del caso de uso.
    \end{cuenum}} \\ \hline
  \textbf{Flujos alternos} & \cuCelda{\textbf{\mbox{FA-01}: El token no corresponde a ninguna verificación}\par
    \begin{cuenum}[start=1]
      \item El Servidor de partidas registra el intento en la bitácora del servidor con la dirección de red de origen, por tratarse de un posible tanteo de enlaces.
      \item El Servidor de partidas muestra la página de error indicando que el enlace no es válido, sin precisar si nunca existió o si ya fue consumido (\mbox{RN-04}).
      \item Fin del caso de uso.
    \end{cuenum}} \\
   & \cuCelda{\textbf{\mbox{FA-02}: El token ya fue usado}\par
    \begin{cuenum}[start=1]
      \item El Servidor de partidas comprueba el \texttt{estado\_\allowbreak{}cuenta} del \textit{Usuario} asociado.
      \item Si la cuenta ya está \texttt{ACTIVA}, el Servidor de partidas muestra una página indicando que la cuenta ya estaba verificada y que puede iniciar sesión. No es un error: reabrir el mismo enlace es lo que hace cualquiera que dude de si funcionó.
      \item Fin del caso de uso.
    \end{cuenum}} \\
   & \cuCelda{\textbf{\mbox{FA-03}: El token expiró}\par
    \begin{cuenum}[start=1]
      \item El Servidor de partidas marca el \textit{TokenVerificacionCorreo} como invalidado.
      \item El Servidor de partidas muestra la página indicando que el enlace caducó y que puede pedir uno nuevo desde el inicio de sesión del juego, que es el \mbox{FA-09} del \mbox{CU-01}.
      \item Fin del caso de uso.
    \end{cuenum}} \\
   & \cuCelda{\textbf{\mbox{FA-04}: La cuenta tiene una sanción de ámbito de cuenta}\par
    \begin{cuenum}[start=1]
      \item El Servidor de partidas no activa la cuenta: verificar el correo no levanta una sanción.
      \item El Servidor de partidas marca el token como usado de todos modos, para que no quede vigente.
      \item El Servidor de partidas muestra la página indicando que la cuenta está sancionada, con el motivo y, si la sanción es temporal, la fecha en que termina.
      \item Fin del caso de uso.
    \end{cuenum}} \\
   & \cuCelda{\textbf{\mbox{FA-05}: La cuenta fue eliminada antes de verificarse}\par
    \begin{cuenum}[start=1]
      \item El Servidor de partidas invalida el token.
      \item El Servidor de partidas muestra la página de enlace no válido, igual que en el \mbox{FA-01}, sin revelar que la cuenta existió (\mbox{RN-04}).
      \item Fin del caso de uso.
    \end{cuenum}} \\
   & \cuCelda{\textbf{\mbox{FA-06}: El Jugador abre el enlace con la sesión ya iniciada en otro dispositivo}\par
    \begin{cuenum}[start=1]
      \item El Servidor de partidas ejecuta el FN sin ninguna variación: la verificación no depende de que haya una sesión abierta ni la afecta.
      \item El SJM del dispositivo con sesión no se entera del cambio hasta su siguiente mensaje al servidor, que ya encontrará la cuenta \texttt{ACTIVA}.
      \item El flujo continúa en el paso 8 del FN.
    \end{cuenum}} \\
   & \cuCelda{\textbf{\mbox{FA-07}: El token confirma un cambio de correo}\par
    \begin{cuenum}[start=1]
      \item El Servidor de partidas encuentra que el \textit{TokenVerificacionCorreo} tiene \texttt{proposito} igual a \texttt{CAMBIO\_\allowbreak{}CORREO} y un \texttt{correo\_\allowbreak{}destino}, y que la cuenta está \texttt{ACTIVA}.
      \item El Servidor de partidas verifica que el \texttt{correo\_\allowbreak{}destino} siga sin estar registrado por otro \textit{Usuario}; si ya lo está, invalida el token y muestra la página de enlace no válido.
      \item El Servidor de partidas sustituye el \texttt{correo} del \textit{Usuario} por el \texttt{correo\_\allowbreak{}destino}, marca el token como usado y registra el cambio en la \textit{BitacoraAcceso} con el resultado \texttt{CAMBIO\_\allowbreak{}CORREO\_\allowbreak{}CONFIRMADO}, en una sola transacción.
      \item El Servidor de partidas muestra la página de confirmación indicando que el correo de la cuenta cambió.
      \item Fin del caso de uso. El \texttt{estado\_\allowbreak{}cuenta} no cambia (\mbox{RN-07} del \mbox{CU-09}).
    \end{cuenum}} \\ \hline
  \textbf{Excepciones} & \cuCelda{\textbf{\mbox{EX-01}: Fallo al consultar o actualizar la base de datos}\par
    \begin{cuenum}[start=1]
      \item El Servidor de partidas revierte la transacción en curso, si la hubiera, de modo que no quede una cuenta activada con el token sin consumir ni al revés.
      \item El Servidor de partidas registra la excepción en la bitácora del servidor con un identificador de incidencia.
      \item El Servidor de partidas muestra la página de error con el identificador, indicando que vuelva a abrir el enlace más tarde. El token sigue vigente, así que el reintento funcionará.
      \item Fin del caso de uso.
    \end{cuenum}} \\
   & \cuCelda{\textbf{\mbox{EX-02}: Fallo al confirmar la transacción}\par
    \begin{cuenum}[start=1]
      \item El Servidor de partidas no obtiene confirmación de que la transacción se haya escrito y la trata como fallida.
      \item El Servidor de partidas registra la excepción señalando que el estado de la activación es indeterminado.
      \item El Servidor de partidas muestra la página de error indicando que intente iniciar sesión antes de volver a abrir el enlace: si la activación sí se escribió, el acceso funcionará y el enlace ya no hará falta.
      \item Fin del caso de uso.
    \end{cuenum}} \\
   & \cuCelda{\textbf{\mbox{EX-03}: La solicitud llega sin token o con un token mal formado}\par
    \begin{cuenum}[start=1]
      \item El Servidor de partidas descarta la solicitud sin consultar la Base de datos.
      \item El Servidor de partidas muestra la página de enlace no válido, idéntica a la del \mbox{FA-01}.
      \item Fin del caso de uso.
    \end{cuenum}} \\ \hline
  \textbf{Postcondiciones} & \cuCelda{\begin{cureglas}
      \item \textbf{\mbox{POST-01}} Si el caso termina por el FN, el \textit{Usuario} tiene \texttt{estado\_\allowbreak{}cuenta} en \texttt{ACTIVA} y una \texttt{fecha\_\allowbreak{}verificacion}, y el \mbox{CU-01} ya no lo rechaza por su \mbox{FA-09}.
      \item \textbf{\mbox{POST-02}} Si el caso termina por el FN, el \textit{TokenVerificacionCorreo} queda marcado como usado y ningún otro token del mismo \textit{Usuario} sigue vigente.
      \item \textbf{\mbox{POST-03}} Si el caso termina por el \mbox{FA-03} o el \mbox{FA-05}, no queda ningún token vigente, pero la cuenta sigue en \texttt{PENDIENTE} y puede pedirse un enlace nuevo.
      \item \textbf{\mbox{POST-04}} Si el caso termina por el \mbox{FA-04}, la cuenta sigue en \texttt{PENDIENTE} y la sanción sigue vigente: ninguna de las dos cosas cambia por haber verificado el correo.
      \item \textbf{\mbox{POST-05}} Si el caso termina por cualquier excepción, ni el \texttt{estado\_\allowbreak{}cuenta} ni el token cambiaron, y el enlace puede volver a abrirse.
    \end{cureglas}} \\ \hline
  \textbf{Reglas de negocio} & \cuCelda{\begin{cureglas}
      \item \textbf{\mbox{RN-01}} El \textit{TokenVerificacionCorreo} es de un solo uso.
      \item \textbf{\mbox{RN-02}} El token tiene una vigencia de veinticuatro horas, la misma que fija el \mbox{RN-11} del \mbox{CU-02}.
      \item \textbf{\mbox{RN-03}} Solo puede haber un token de verificación vigente por cuenta. Consumir uno invalida el resto.
      \item \textbf{\mbox{RN-04}} Un enlace inexistente, uno de una cuenta eliminada y uno manipulado producen la misma página de error. El proceso no puede usarse para averiguar qué cuentas existen.
      \item \textbf{\mbox{RN-05}} Verificar el correo cambia únicamente el \texttt{estado\_\allowbreak{}cuenta}. No levanta sanciones, no abre sesión y no otorga nada.
      \item \textbf{\mbox{RN-06}} El token se almacena como resumen criptográfico, nunca en claro, por el mismo motivo que la contraseña.
      \item \textbf{\mbox{RN-07}} Este es el único flujo del sistema que puede llevar una cuenta de \texttt{PENDIENTE} a \texttt{ACTIVA}.
    \end{cureglas}} \\ \hline
  \textbf{Observación} & \cuCelda{\textit{Sobre por qué este caso no abre sesión.}\par
    Sería cómodo que verificar el correo dejara al Jugador dentro del juego, pero el enlace se abre en un navegador y la sesión vive en el cliente de escritorio: no hay forma de pasar el token de sesión de uno a otro sin inventar un mecanismo que no existe en el resto del sistema. El Jugador verifica, vuelve al juego e inicia sesión con el \mbox{CU-01}, que es el único lugar donde se crea una \textit{Sesion}.} \\ \hline
  \textbf{Datos que toca} & \cuCelda{\begin{cudatos}
      \item \textbf{\textit{TokenVerificacionCorreo}:} lee \texttt{token\_\allowbreak{}hash}, \texttt{usado}, \texttt{fecha\_\allowbreak{}expiracion} \textit{(pasos 3, 5, 6)}
      \item \textbf{\textit{TokenVerificacionCorreo}:} marca como usado e invalida los demás \textit{(pasos 10, 11)}
      \item \textbf{\textit{Usuario}:} lee \texttt{estado\_\allowbreak{}cuenta}, \texttt{idioma\_\allowbreak{}preferido} \textit{(pasos 7, 14)}
      \item \textbf{\textit{Usuario}:} escribe \texttt{estado\_\allowbreak{}cuenta}, \texttt{fecha\_\allowbreak{}verificacion} \textit{(paso 9)}
      \item \textbf{\textit{Sancion}:} lee, para el \mbox{FA-04} \textit{(paso 7)}
      \item \textbf{\textit{BitacoraAcceso}:} inserta \textit{(paso 13)}
    \end{cudatos}} \\ \hline
\end{fichacu}
\clearpage
